% Volume IV: Optimization and Computation
% Continuity note for Chapter 24.
% Previous dependencies: Chapters 4-6 provide linear systems, rank, projections, least squares, QR, eigendecomposition, SVD, and quadratic forms; Chapter 7 provides matrix calculus and tensor contractions; Chapters 20-23 provide optimization objectives, gradients, computational graphs, floating-point stability, and vectorized computation.
% New durable concepts: conditioning; ill-posed inverse problems; singular-value sensitivity; stable solves instead of explicit inverses; LU, Cholesky, QR, eigendecomposition, and SVD in computation; direct versus iterative solvers; residuals; Krylov methods; conjugate gradient; preconditioning; matrix-free products; discretization; local and global truncation error; stability, consistency, convergence, and scientific-computing workflows.
% Forward hooks: Chapter 25 uses numerical integration for ODEs; Chapter 26 uses discretization for PDEs and continuum models; Chapter 31 onward reuse least-squares, regularization, large optimization, kernels, Gaussian processes, and differentiable simulation; neuroscience chapters use stable simulation and inverse problems for neural data analysis.
% Notation commitments: Linear systems are Ax=b with A in R^{m x n}; square solves usually use A in R^{n x n}; perturbations are \Delta A,\Delta b,\Delta x; matrix condition numbers use \kappa_2(A)=\sigma_{\max}(A)/\sigma_{\min}(A) for full-rank matrices; residuals are r_k=b-Ax_k; discretization step sizes are \Delta t and h.
\section{Chapter 24: Numerical Linear Algebra and Scientific Computing}
\label{ch:numerical-linear-algebra-scientific-computing}

Chapter 23 showed how programs compute derivatives. This chapter asks whether the underlying numerical computations are trustworthy. A formula can be mathematically correct and still be a poor way to compute. Solving a linear system by explicitly forming \(A^{-1}\), squaring a condition number through normal equations, or taking a time step outside a stability region can turn small measurement noise into large numerical artifacts.

Numerical linear algebra is the bridge between finite-dimensional mathematics and actual scientific computation. It studies sensitivity of problems, stability of algorithms, and the structure that lets large computations be done without wasting memory or accuracy.

\paragraph{Chapter problem.}
How do we solve matrix problems, fit models, and simulate systems in ways that respect conditioning, floating-point arithmetic, scale, and discretization error?

\paragraph{Section map.}
Section 24.1 defines conditioning and ill-posedness through perturbations and singular values. Section 24.2 explains the main matrix factorizations used in computation. Section 24.3 studies direct and iterative methods for large systems. Section 24.4 introduces simulation, discretization, truncation error, and stability.

\subsection{24.1 Conditioning and ill-posedness}

\paragraph{Guiding question.}
When does a small error in data or measurements cause a large error in the computed answer?

\paragraph{Beginner intuition.}
Conditioning is about the problem, not the algorithm. A well-conditioned problem has outputs that change only moderately when inputs are slightly perturbed. An ill-conditioned problem has outputs that can change dramatically under tiny perturbations. In data analysis, ill-conditioning often means that two explanatory directions are nearly redundant or that the measurements do not contain enough information to recover the desired unknown.

The smallest example is a diagonal linear system
\[
A x=b,\qquad
A=\begin{bmatrix}1&0\\0&0.01\end{bmatrix}.
\]
If \(b=(1,0.01)^\top\), then \(x=(1,1)^\top\). But if the second component of \(b\) is perturbed from \(0.01\) to \(0.011\), the second component of \(x\) changes from \(1\) to \(1.1\). A tiny absolute measurement change has been magnified by the small singular direction of \(A\).

\paragraph{Formal setup.}
For a scalar function \(f\), a local relative condition number can be written informally as
\[
\operatorname{cond}(f,x)
\approx
\frac{\text{relative output change}}{\text{relative input change}}
\]
for small perturbations. For solving a full-rank square system
\[
Ax=b,\qquad A\in\mathbb{R}^{n\times n},
\]
the Euclidean matrix condition number is
\[
\boxed{\kappa_2(A)=\|A\|_2\|A^{-1}\|_2
=\frac{\sigma_{\max}(A)}{\sigma_{\min}(A)}.}
\]
Here \(\sigma_{\max}\) and \(\sigma_{\min}\) are the largest and smallest singular values. If \(A\) is singular, then \(\sigma_{\min}=0\) and \(\kappa_2(A)=\infty\).

When only \(b\) is perturbed and \(A\) is fixed, the solution perturbation satisfies the standard bound
\[
\frac{\|\Delta x\|_2}{\|x\|_2}
\leq
\kappa_2(A)
\frac{\|\Delta b\|_2}{\|b\|_2},
\]
assuming \(A\) is invertible and \(x,b\neq 0\). This is a worst-case upper bound; it says how much amplification the matrix can allow.

An inverse problem is \emph{ill-posed} when the desired cause is not stably determined by the observations. In finite dimensions, tiny singular values are the signature: directions in the unknown are barely visible in the data.

\paragraph{Worked mechanism: singular values and regularization.}
Let \(A=U\Sigma V^\top\) be the SVD of a full-rank matrix, with singular values \(\sigma_i>0\). Solving \(Ax=b\) gives
\[
x=V\Sigma^{-1}U^\top b.
\]
In coordinates,
\[
x=\sum_i \frac{u_i^\top b}{\sigma_i}v_i.
\]
If \(\sigma_i\) is tiny, then noise in the data coefficient \(u_i^\top b\) is divided by a tiny number. This amplifies noise.

Ridge or Tikhonov regularization replaces the unstable least-squares problem by
\[
\min_x \frac12\|Ax-b\|_2^2+\frac{\lambda}{2}\|x\|_2^2,
\qquad \lambda>0.
\]
The normal equation becomes
\[
(A^\top A+\lambda I)x=A^\top b.
\]
In SVD coordinates,
\[
x_\lambda=\sum_i
\frac{\sigma_i}{\sigma_i^2+\lambda}(u_i^\top b)v_i.
\]
The factor \(\sigma_i/(\sigma_i^2+\lambda)\) does not explode as violently when \(\sigma_i\) is small. Regularization trades some bias for reduced variance and improved stability.

\paragraph{Examples and connections.}
Numerically, for \(A=\operatorname{diag}(1,0.01)\),
\[
\kappa_2(A)=100.
\]
The second coordinate is sensitive because the second singular value is \(0.01\). If \(b_2\) changes by \(10^{-3}\), then \(x_2=b_2/0.01\) changes by \(0.1\).

In computational neuroscience, estimating a stimulus from neural responses can be ill-conditioned if many neurons have nearly identical tuning curves. The response matrix has nearly redundant columns, so different stimuli or latent variables produce almost the same population response. Regularization encodes the decision not to trust barely observed directions.

In AI, ill-conditioning appears in least-squares layers, kernel regression, Gaussian-process covariance matrices, and deep-network curvature. Optimization may be slow in directions where curvature is tiny and unstable in directions where curvature is huge.

\paragraph{Common misconceptions and failure modes.}
Conditioning is not probability conditioning. It is numerical sensitivity. A second misconception is that a small residual always means an accurate solution. For ill-conditioned systems, a small residual can coexist with a large solution error. A third failure is to compute \(A^{-1}b\) explicitly. The inverse can be more expensive and less stable than solving the system through a factorization.

\paragraph{Exercises.}
\begin{enumerate}[label=\textbf{24.1.\arabic*.}, leftmargin=*]
\item \textbf{Mechanical.} Compute \(\kappa_2(A)\) for \(A=\operatorname{diag}(2,0.5)\).
\item \textbf{Proof or derivation.} Using \(x=A^{-1}b\) and \(\Delta x=A^{-1}\Delta b\), derive the perturbation bound \(\|\Delta x\|_2/\|x\|_2\leq \kappa_2(A)\|\Delta b\|_2/\|b\|_2\).
\item \textbf{Numerical/computational.} For \(A=\operatorname{diag}(1,0.001)\), \(b=(1,0.001)^\top\), compute \(x\). Then change \(b_2\) to \(0.0012\) and compute the new solution.
\item \textbf{Interpretation.} In a neural decoding matrix with two almost identical columns, why is the inverse problem of separating the two corresponding features unstable?
\item \textbf{Misconception/debugging.} A student says, "My residual \(\|Ax-b\|\) is small, so my solution must be accurate." Explain why this can fail for ill-conditioned \(A\).
\end{enumerate}

\subsection{24.2 Matrix factorizations in computation}

\paragraph{Guiding question.}
Which matrix decomposition should we use for a numerical task, and why is it usually better than forming an inverse?

\paragraph{Beginner intuition.}
A factorization rewrites a matrix as a product of structured matrices. The point is not only theoretical understanding. The structure makes computation stable, efficient, and interpretable. Triangular systems are easy to solve. Orthogonal matrices preserve lengths. Diagonal singular-value factors reveal rank and sensitivity.

The main habit is: solve systems by factoring \(A\), not by computing \(A^{-1}\). If \(Ax=b\), a good numerical library usually computes a factorization and then solves simpler triangular or orthogonal subproblems.

\paragraph{Formal setup.}
Common factorizations include:
\[
A=LU
\]
for square systems, usually with row pivoting \(PA=LU\);
\[
A=QR,
\]
where \(Q^\top Q=I\) and \(R\) is upper triangular;
\[
A=LL^\top
\]
for symmetric positive definite matrices, called Cholesky factorization;
\[
A=V\Lambda V^{-1}
\]
for diagonalizable square matrices; and
\[
A=U\Sigma V^\top
\]
for the singular value decomposition.

\begin{center}
\small
\begin{tabular}{@{}p{0.18\linewidth}p{0.22\linewidth}p{0.24\linewidth}p{0.25\linewidth}@{}}
\toprule
Factorization & Main assumption & Typical use & Stability idea\\
\midrule
LU with pivoting & square, nonsingular & general solves & pivot away from tiny pivots\\
Cholesky & symmetric positive definite & covariance and normal systems & exploit positivity\\
QR & full-column-rank or rectangular & least squares & orthogonal transformations\\
Eigen & square with spectral structure & dynamics, modes & invariant directions\\
SVD & any matrix & rank, conditioning, pseudoinverse & singular directions\\
\bottomrule
\end{tabular}
\end{center}

\paragraph{Worked mechanism: QR least squares.}
For an overdetermined least-squares problem
\[
\min_x \|Ax-b\|_2^2,\qquad A\in\mathbb{R}^{m\times n},\quad m\geq n,
\]
suppose \(A\) has full column rank and thin QR factorization
\[
A=QR,
\qquad
Q\in\mathbb{R}^{m\times n},\quad Q^\top Q=I,\quad R\in\mathbb{R}^{n\times n}.
\]
Then
\[
\|Ax-b\|_2^2=\|QRx-b\|_2^2.
\]
Orthogonal projection separates \(b\) into a part in the column space of \(Q\) and an orthogonal residual. The minimizer satisfies
\[
Rx=Q^\top b.
\]
Since \(R\) is triangular, solve this by back substitution.

This avoids the normal equations
\[
A^\top A x=A^\top b,
\]
which square the condition number:
\[
\kappa_2(A^\top A)=\kappa_2(A)^2
\]
when \(A\) has full column rank. Squaring a large condition number can destroy accuracy.

\paragraph{Examples and connections.}
Numerically, let
\[
A=\begin{bmatrix}1&0\\1&1\\1&2\end{bmatrix},\qquad b=\begin{bmatrix}1\\2\\2\end{bmatrix}.
\]
Least squares fits a line \(x_1+x_2 t\) to the three data points. QR gives an orthonormal basis for the two-dimensional column space and solves a \(2\times 2\) triangular system. The residual is the part of \(b\) orthogonal to that column space.

In computational neuroscience, SVD and QR appear when fitting low-dimensional latent trajectories to high-dimensional neural recordings. SVD identifies dominant population activity patterns, while QR gives stable least-squares decoders. In AI, factorizations appear in linear layers, whitening, second-order optimization, Gaussian-process covariance solves, and low-rank adaptation.

\paragraph{Common misconceptions and failure modes.}
The inverse \(A^{-1}\) is a mathematical object, not a recommended computational plan. Another failure is to use eigendecomposition for every matrix; nonnormal matrices can have misleading eigenvectors, and rectangular matrices need SVD instead. A third misconception is that normal equations are always harmless because they are algebraically correct. They may be acceptable for well-conditioned small problems, but QR or SVD is safer when accuracy matters.

\paragraph{Exercises.}
\begin{enumerate}[label=\textbf{24.2.\arabic*.}, leftmargin=*]
\item \textbf{Mechanical.} Match each task to a factorization: solving an SPD covariance system, least-squares fitting, estimating numerical rank, and simulating modes of a diagonalizable linear system.
\item \textbf{Proof or derivation.} Show that if \(A=QR\) with \(Q^\top Q=I\), then the least-squares normal equation \(A^\top(Ax-b)=0\) becomes \(R^\top(Rx-Q^\top b)=0\). If \(R\) is invertible, conclude \(Rx=Q^\top b\).
\item \textbf{Numerical/computational.} If \(\kappa_2(A)=10^4\), what is \(\kappa_2(A^\top A)\) for full-column-rank \(A\)? Why is this a warning about normal equations?
\item \textbf{Interpretation.} Explain why SVD is useful for identifying low-dimensional neural population structure.
\item \textbf{Misconception/debugging.} A student computes \(x=A^{-1}b\) inside a loop for many right-hand sides. Explain a better factorization-based workflow.
\end{enumerate}

\subsection{24.3 Solving large systems}

\paragraph{Guiding question.}
How can we solve \(Ax=b\) when \(A\) is too large to factor or even store densely?

\paragraph{Beginner intuition.}
Large scientific systems are often sparse or structured. A finite-difference PDE, a graph Laplacian, or a recurrent network linearization may involve millions of variables, but each variable interacts directly with only a few neighbors. Direct dense factorization can be impossible. Iterative methods instead build a sequence of approximations \(x_k\) using matrix-vector products, residuals, and sometimes a preconditioner.

\paragraph{Formal setup.}
For a candidate solution \(x_k\), the residual is
\[
r_k=b-Ax_k.
\]
It measures how far \(x_k\) is from satisfying the equation. Iterative solvers stop when a norm such as
\[
\frac{\|r_k\|_2}{\|b\|_2}
\]
is below a chosen tolerance, while also checking whether the problem is ill-conditioned enough that a small residual may still be misleading.

A method is \emph{matrix-free} if it needs only products \(v\mapsto Av\), not explicit storage of all entries of \(A\). Krylov methods build approximations in spaces such as
\[
\mathcal{K}_k(A,r_0)=
\operatorname{span}\{r_0,Ar_0,A^2r_0,\ldots,A^{k-1}r_0\}.
\]
For symmetric positive definite \(A\), conjugate gradient is the standard Krylov method.

A preconditioner \(M\) is an approximate easy-to-solve version of \(A\). Instead of solving \(Ax=b\) directly, one solves a transformed problem such as
\[
M^{-1}Ax=M^{-1}b.
\]
The goal is to reduce the effective condition number while keeping applications of \(M^{-1}\) cheap.

\paragraph{Worked mechanism: conjugate-gradient viewpoint.}
If \(A=A^\top\) is positive definite, solving \(Ax=b\) is equivalent to minimizing the quadratic
\[
\phi(x)=\frac12x^\top Ax-b^\top x.
\]
The gradient is
\[
\nabla\phi(x)=Ax-b=-r.
\]
Conjugate gradient chooses search directions that are conjugate with respect to \(A\), meaning
\[
p_i^\top A p_j=0\quad\text{for }i\neq j.
\]
This avoids repeatedly correcting the same error direction. In exact arithmetic, conjugate gradient solves an \(n\)-dimensional SPD problem in at most \(n\) steps, but its practical power is that a good approximate solution often appears much earlier when the spectrum is favorable or preconditioned.

\paragraph{Examples and connections.}
Numerically, suppose \(A=\begin{bmatrix}4&1\\1&3\end{bmatrix}\), \(b=(1,2)^\top\), and \(x_0=(0,0)^\top\). The first residual is \(r_0=b\). A steepest descent step along \(r_0\) uses
\[
\alpha_0=\frac{r_0^\top r_0}{r_0^\top A r_0}
=\frac{5}{(1,2)\cdot(6,7)}
=\frac{5}{20}
=0.25,
\]
so \(x_1=(0.25,0.5)^\top\). Conjugate gradient improves on steepest descent by choosing later directions that account for the geometry of \(A\).

In computational neuroscience, implicit simulations of cable equations or network diffusion models produce large sparse linear systems at each time step. Iterative solvers exploit sparsity and local connectivity. In AI, large systems arise in kernel methods, Gaussian-process inference, second-order optimization, natural gradients, and implicit layers. Hessian-vector products from automatic differentiation allow matrix-free curvature methods without materializing the Hessian.

\paragraph{Common misconceptions and failure modes.}
A small residual tolerance is not a universal accuracy guarantee; conditioning still matters. Iterative methods are not automatically faster than direct methods; for small dense systems, a direct factorization may be better. Preconditioning is not optional decoration for difficult problems; it can determine whether an iterative method converges in useful time. Finally, matrix-free does not mean cost-free: each matrix-vector product may still be expensive.

\paragraph{Exercises.}
\begin{enumerate}[label=\textbf{24.3.\arabic*.}, leftmargin=*]
\item \textbf{Mechanical.} For \(A=\begin{bmatrix}2&0\\0&5\end{bmatrix}\), \(b=(2,10)^\top\), and \(x=(1,1)^\top\), compute the residual.
\item \textbf{Proof or derivation.} Show that the gradient of \(\phi(x)=\frac12x^\top Ax-b^\top x\), with \(A=A^\top\), is \(Ax-b\). Conclude that the minimizer satisfies \(Ax=b\).
\item \textbf{Numerical/computational.} Compute the first steepest-descent step for \(A=\begin{bmatrix}3&0\\0&1\end{bmatrix}\), \(b=(3,1)^\top\), and \(x_0=(0,0)^\top\) using \(\alpha_0=(r_0^\top r_0)/(r_0^\top A r_0)\).
\item \textbf{Interpretation.} Why is sparsity useful when simulating a spatial neural model with only local interactions?
\item \textbf{Misconception/debugging.} A student forms a dense \(10^6\times 10^6\) matrix for a local-neighbor PDE discretization. Explain why a matrix-free or sparse representation is usually required.
\end{enumerate}

\subsection{24.4 Simulation, discretization, and error}

\paragraph{Guiding question.}
When a continuous mathematical model is simulated on a finite computer, where does error enter and how can we control it?

\paragraph{Beginner intuition.}
Scientific computing turns continuous equations into finite computations. Time is replaced by time steps, space by grids or basis functions, and differential operators by finite-dimensional approximations. This process is called discretization. The result is never exact; it has truncation error from replacing a continuous equation by a discrete rule, roundoff error from floating-point arithmetic, and often modeling and statistical error from the data-generating assumptions.

\paragraph{Formal setup.}
For an ODE
\[
\dot x(t)=f(x(t),t),\qquad x(t_0)=x_0,
\]
the forward Euler discretization with time step \(\Delta t\) is
\[
x_{k+1}=x_k+\Delta t\, f(x_k,t_k),
\qquad t_k=t_0+k\Delta t.
\]
The local truncation error is the one-step error made when the exact solution is inserted into the numerical update. The global error is the accumulated error after many steps.

A numerical method is \emph{consistent} if its local approximation error goes to zero as \(\Delta t\to 0\). It is \emph{stable} if errors are not amplified uncontrollably by the time-stepping rule. It is \emph{convergent} if the numerical solution approaches the true solution as the discretization is refined. In many standard settings, consistency plus stability implies convergence.

\paragraph{Worked mechanism: Euler stability for a linear test equation.}
Consider
\[
\dot x=\lambda x,\qquad \lambda<0.
\]
The exact solution is
\[
x(t)=e^{\lambda t}x(0),
\]
which decays to zero. Forward Euler gives
\[
x_{k+1}=x_k+\Delta t\,\lambda x_k=(1+\Delta t\,\lambda)x_k.
\]
Therefore
\[
x_k=(1+\Delta t\,\lambda)^k x_0.
\]
For the numerical solution to decay, we need
\[
|1+\Delta t\,\lambda|<1.
\]
If \(\lambda=-10\), then \(\Delta t=0.1\) gives multiplier \(0\), which is stable but crude. The step \(\Delta t=0.3\) gives multiplier \(-2\), so the numerical solution oscillates with growing magnitude even though the true solution decays. The problem is not the ODE; it is the time step and method.

\paragraph{Examples and connections.}
A numerical simulation of the membrane equation
\[
\tau \dot V=-(V-E_L)+RI(t)
\]
uses a time step \(\Delta t\). If \(\Delta t\) is too large compared with \(\tau\), explicit Euler can create artificial oscillations or blow-up. Smaller steps, implicit methods, or exact updates for linear subproblems may be needed.

For PDEs, space is also discretized. A one-dimensional second derivative may be approximated by
\[
\frac{\partial^2 u}{\partial x^2}(x_i)
\approx
\frac{u_{i+1}-2u_i+u_{i-1}}{h^2}.
\]
This converts a continuum model into a coupled finite-dimensional system. Chapter 26 will use this idea for diffusion, waves, and neural fields.

In AI, neural ODEs and differentiable physics models depend on numerical solvers. The learned function \(f_\theta\) is only part of the model; solver tolerance, step-size control, stiffness, and adjoint accuracy also affect predictions and gradients.

\paragraph{Common misconceptions and failure modes.}
More steps are not always better if roundoff or computational cost dominates, although refinement usually reduces truncation error in a stable method. A second misconception is that a visually smooth trajectory is necessarily accurate. Numerical damping can hide wrong dynamics. A third failure is to compare models without fixing solver tolerances; the apparent model difference may be a simulation artifact.

\paragraph{Exercises.}
\begin{enumerate}[label=\textbf{24.4.\arabic*.}, leftmargin=*]
\item \textbf{Mechanical.} Apply one forward Euler step to \(\dot x=-2x\), \(x_0=3\), with \(\Delta t=0.1\).
\item \textbf{Proof or derivation.} Derive the Euler stability condition \(|1+\Delta t\,\lambda|<1\) for \(\dot x=\lambda x\).
\item \textbf{Numerical/computational.} For \(\dot x=-5x\), compare the Euler multipliers for \(\Delta t=0.1\), \(0.3\), and \(0.5\). Which are stable?
\item \textbf{Interpretation.} In a conductance-based neuron simulation, why should the time step be chosen relative to the fastest relevant time constant?
\item \textbf{Misconception/debugging.} A simulation of a decaying system grows in magnitude after switching to a larger \(\Delta t\). Explain how numerical instability can cause this even when the continuous model is stable.
\end{enumerate}

\paragraph{Closing bridge.}
This chapter turned linear algebra and calculus into practical computation: sensitivity, stable factorizations, scalable solves, and discretization. The next chapter studies the continuous-time systems that many of these computations simulate. There, the same themes reappear as flows, equilibria, stability, phase portraits, and numerical integration of ordinary differential equations.

\clearpage
